\section{Introduction}
Financial institutions ingest checks, identification forms, and statements as images, and downstream systems need to know which string is the account number, which is the date, and which is the amount. Page-level OCR returns text but not that mapping, and its accuracy is not a fixed number: published evaluations report Tesseract anywhere from 47 percent accuracy on Thai vehicle registration certificates \cite{thammarak2022} to 92.6 percent word recognition on letters \cite{reisswig2019}, with different metrics, document types, resolutions, and preprocessing. A detect-then-read design addresses both problems. An object detector localizes each field, and OCR runs only on the cropped region, so every extracted string arrives already labeled with its field class.

Ultralytics YOLO26 \cite{jocher2026yolo26} is a natural detector for this role. Its end-to-end head does not need non-maximum suppression at inference, which removes a post-processing pass from the serving path. Running it over millions of documents calls for a distributed engine, and Apache Spark \cite{zaharia2012} is the usual choice.

The difficulty is that a model call inside a Spark job looks routine and fails quietly. The RAPIDS Accelerator for Apache Spark \cite{nvidia_udf,nvidia_rapids} accelerates Spark\textquotesingle{}s own SQL and DataFrame plan on GPU. It does not accelerate Python code that loads a model inside a user-defined function, so enabling it does not make model inference fast. The first version of this pipeline contained several such defects, some in the model call pattern, some in data access, and some in the toolchain. We found them only after running the code, which is the reason this paper reports a running system and not a design alone.

This paper builds the pipeline, runs it, measures it, and releases it. Its scope is deliberately narrow. All reported accuracy is on a synthetic form dataset whose labels are exact by construction, timing is on a two-core CPU with one replication on a single GPU, and the AWS EMR and multi-GPU deployment is presented as design, not as measurement.

\section{Contributions}
\begin{itemize}
\item A complete open implementation of field detection, field-level OCR, and aggregation as one Spark job, with unit and end-to-end tests. One command reproduces the CPU results, and a notebook repeats the study on a GPU.
\item A catalog of integration defects, marked by whether we reproduced them by running code or found them by inspection, and the corrected partition-level inference pattern that resolves them (Section~\ref{sec:impl}).
\item An evaluation on a dataset with exact ground truth, with one training run per device: detection precision and recall, end-to-end field accuracy against a full-page Tesseract baseline, and a measured comparison of three ways to run the same model inside Spark (Section~\ref{sec:results}).
\item A deployment design for AWS EMR with RAPIDS, kept separate from what was measured, including when the batch design is and is not appropriate (Section~\ref{sec:design}).
\end{itemize}

\section{Related Work}
\textbf{Detection.} Single-network detection was introduced by Redmon et al.\ \cite{redmon2016}. YOLOv9 added programmable gradient information \cite{wang2024yolov9}, and YOLOv10 was the first in the line to remove non-maximum suppression \cite{wang2024yolov10}. YOLO26 keeps native end-to-end inference, removes Distribution Focal Loss, and is a recent Ultralytics release \cite{jocher2026yolo26}. DocLayout-YOLO adapts the family to document layout analysis \cite{zhao2024doclayout}. We use YOLO26 because it is the maintained production model and its NMS-free head simplifies the serving path. We do not claim it is the best detector for financial fields, and we do not compare against other detectors.

\textbf{OCR and document understanding.} Tesseract \cite{smith2007} remains a common open-source engine, and its accuracy varies widely with document type and preprocessing \cite{thammarak2022,reisswig2019}. Managed services such as Amazon Textract \cite{textract} are an obvious comparison, which we do not make because it needs an AWS account and real documents. Layout-aware transformers such as LayoutLMv3 \cite{huang2022layoutlmv3} and OCR-free models such as Donut \cite{kim2022donut} are stronger, heavier alternatives that solve a broader problem. Our baseline is Tesseract alone, which makes the comparison one of pipeline structure, not one against the state of the art.

\textbf{Datasets.} FUNSD \cite{jaume2019} and SROIE \cite{huang2019sroie} are the public datasets closest to this domain, and neither fits. FUNSD annotates form entities as question, answer, header, and other. SROIE provides text-line boxes and key-information strings for receipts. Neither provides region boxes for classes such as account number, date, amount, name, and signature. For that reason we evaluate on a synthetic dataset whose boxes and strings are exact, and we say plainly what that does and does not show.

\textbf{Systems.} Spark\textquotesingle{}s resilient distributed datasets \cite{zaharia2012} underlie its DataFrame engine. Spark exchanges data with Python through Apache Arrow, and \texttt{mapInPandas} \cite{spark_mapinpandas} exposes that path at partition granularity. Spark can schedule GPUs as task resources \cite{spark_resources}, and the RAPIDS documentation describes how GPU-scheduled Python workers share a device with the executor \cite{nvidia_udf}.
